\chapter{Introduction}
\label{ch:introduction}

\section{Background and Motivation}

Biometric authentication is now deeply embedded in critical infrastructure.
Border control systems, national identity programs, mobile banking, and access
management platforms routinely rely on fingerprint, iris, and face recognition to
verify who a person is \cite{jain2004multibiometrics}. Unlike passwords or tokens,
biometric traits are permanent: a user cannot change their fingerprint the way they
change a password. This permanence creates an acute problem when biometric
\emph{templates}---the stored representations used for matching---are compromised.

ISO/IEC~24745 \cite{iso24745} codifies three mandatory properties for any
\emph{Biometric Template Protection} (BTP) scheme. A protected template must be
\textbf{irreversible}: it must not be possible to reconstruct the original
biometric trait from the stored data. It must be \textbf{unlinkable}: templates
enrolled at different systems (using different transformation parameters) must not
be linkable back to the same individual. And it must be \textbf{revocable}: when a
template is compromised, it must be possible to re-enroll with a fresh, unlinkable
transform. These three requirements together define the security floor for any
system claiming to protect biometric privacy.

The \emph{fuzzy vault} \cite{juels2002fuzzy} is one of the most extensively studied
BTP primitives. It hides a secret polynomial among a large set of chaff points;
only a biometric query that matches the enrolled data sufficiently closely can
identify enough genuine points to reconstruct the polynomial by Lagrange
interpolation \cite{rathgeb2011survey}. Maurya et al.\ \cite{maurya2025enhancing}
advanced the state of the art by encoding quantized minutiae as elliptic curve (EC)
points over Brainpool~P256r1 \cite{rfc5639}, achieving an Equal Error Rate (EER)
as low as 0.005\% on FVC2002~DB1 with vault entropy exceeding 7.7~bits per
coordinate. Despite this progress, unimodal fingerprint systems remain vulnerable
to spoof attacks and performance degrades significantly on challenging data: EER
reaches 5.0\% on FVC2002~DB3.

Multimodal biometric fusion---combining evidence from fingerprint, iris, face, or
other traits---is the standard approach to improving accuracy and spoof
resistance \cite{ross2003information}. However, integrating multimodal fusion
\emph{inside} a template protection scheme introduces a fundamental conflict:
the cryptographic binding that protects the fingerprint vault must simultaneously
accommodate iris data without compromising either modality's security or the
scheme's ISO/IEC~24745 compliance.

\subsection{The Gating Penalty: A Structural Problem}

Every existing multimodal fuzzy vault scheme solves this integration problem by
making the auxiliary modality a \emph{mandatory gate}: both modalities must
succeed for the vault to be unlocked \cite{nandakumar2008multibiometric,
moon2010multibiometric}. The consequence is immediate. Under AND-fusion, the
Genuine Acceptance Rate (GAR) of the combined system is:
\begin{equation}
  \mathrm{GAR}_{\mathrm{AND}} = \mathrm{GAR}_{\mathrm{FP}}
  \times \mathrm{GAR}_{\mathrm{Iris}} \leq \mathrm{GAR}_{\mathrm{FP}}
\end{equation}
Since $\mathrm{GAR}_{\mathrm{Iris}} < 1$ in any realistic deployment, the
combined system is always worse than the fingerprint alone for genuine users.
With the iris GAR measured at approximately 62\% in this work, over one third
of genuine users are rejected by the combined system even when their fingerprint
would have accepted them.

This is not a calibration problem or a dataset artefact. It is a structural
consequence of the AND-fusion architecture: using the iris as a mandatory gate
means that iris failure becomes user rejection, regardless of fingerprint quality.
The same structural flaw affects polynomial blinding (Architecture B), dense chaff
injection (Architecture C), and AES-256-GCM vault encryption (Architecture D)---
all tight-coupling strategies examined in this thesis.

\subsection{The Centralization Problem}

Even if the biometric performance problem were solved, a separate infrastructure
problem remains. Centralized vault storage is a single point of failure. The
Aadhaar breach of 2018 exposed biometric records of over 1.1 billion individuals;
the U.S.\ Office of Personnel Management breach of 2015 exposed 5.6 million
fingerprint records. Once a vault is leaked, an attacker can mount offline
brute-force attacks without time pressure.

Threshold secret sharing---specifically Shamir's $(t,n)$-scheme
\cite{shamir1979how}---offers an information-theoretic solution: split the vault
across $n$ nodes such that any $t$ shares reconstruct it exactly but any
$t{-}1$ or fewer shares reveal zero information. No prior work has applied
vault-level secret sharing to a fuzzy vault, nor combined it with a zero-penalty
multimodal fusion scheme.

\section{Problem Statement}

This thesis addresses two tightly linked problems:

\begin{enumerate}
  \item \textbf{The gating penalty.} All existing multimodal fuzzy vault schemes
  gate decoding on the auxiliary modality. When the iris fails, the user is
  rejected even if the fingerprint alone would have succeeded. A scheme is needed
  where a second modality can help when available but never hurt when unavailable.

  \item \textbf{Centralized vault storage.} Storing protected biometric vaults on
  a single server creates a single point of failure. A distributed storage
  architecture is needed that provides information-theoretic security,
  Byzantine fault tolerance, and auditability.
\end{enumerate}

\section{Objectives}

\begin{enumerate}
  \item To design a multimodal BTP scheme that uses the iris as an optional
  booster rather than a mandatory gate, and to formally prove $\mathrm{EER_{IBFV}}
  \leq \mathrm{EER_{Uni}}$ for all configurations.

  \item To implement and evaluate the proposed IBFV scheme on four FVC fingerprint
  databases (FVC2002 DB1--DB3, FVC2004 DB1) paired with CASIA-Iris-Interval
  chimeric samples, across 960+ parameter configurations.

  \item To compare IBFV against four alternative multimodal architectures
  (AND-fusion, polynomial blinding, dense chaff injection, AES-256-GCM
  encryption) and characterize why tight coupling structurally degrades performance.

  \item To design D-IBFV, a decentralized extension applying $(t,n)$-Shamir Secret
  Sharing to the vault itself, distributing shares on IPFS, and anchoring CIDs on
  a blockchain, with formal proofs of share confidentiality and exact reconstruction.

  \item To validate D-IBFV on Ethereum-compatible (Ganache) and Hyperledger
  Fabric~2.5 platforms, measuring enrollment latency, verification latency,
  throughput, storage costs, node failure resilience, and accuracy.

  \item To demonstrate full ISO/IEC~24745 compliance: irreversibility (ECDLP
  $\sim\!2^{128}$), unlinkability ($D_{\mathrm{sys}} \leq 0.03$), and
  revocability via parameter-based re-enrollment.
\end{enumerate}

\section{Contributions}

\begin{enumerate}
  \item \textbf{IBFV: The first zero-penalty multimodal fuzzy vault.} The
  Iris-Boosted Fuzzy Vault uses the iris as an optional booster by injecting $K$
  iris-derived bonus polynomial evaluation points into the fingerprint vault.
  The zero-penalty guarantee is formally proven (Theorem~1) and empirically
  validated across 960+ configurations. All four alternative architectures degrade
  performance, confirming gating is a structural---not incidental---problem.

  \item \textbf{Full ISO/IEC~24745 security analysis.} A complete analysis against
  Ratha et al.'s eight-point threat model \cite{ratha2001enhancing}: irreversibility
  (ECDLP $\sim\!2^{128}$, iris code recovery $2^{514}$), unlinkability
  ($D_{\mathrm{sys}} \leq 0.03$), and revocability.

  \item \textbf{D-IBFV: Threshold-distributed vault storage.} The first scheme to
  split a fuzzy vault via $(t,n)$-Shamir Secret Sharing over $\mathrm{GF}(257)$,
  with formal proofs of share confidentiality and exact reconstruction, and
  blockchain-anchored integrity.

  \item \textbf{Multi-platform system benchmarks.} Enrollment under 806\,ms,
  verification under 71\,ms, zero-bit-error reconstruction across 30{,}000 trials,
  O(1) verification scaling, and only 0.5\,KB on-chain per user.
\end{enumerate}

\section{Thesis Organization}

Chapter~\ref{ch:related} surveys related work. Chapter~\ref{ch:methodology}
presents the proposed architecture in full technical detail. Chapter~\ref{ch:results}
reports biometric and system experiments. Chapter~\ref{ch:discussion} discusses
implications, security, and limitations. Chapter~\ref{ch:conclusions} concludes
with future directions. Appendix~A provides mathematical proofs.
Appendix~B gives complete algorithm listings.

